% Generated draft: requires mathematical and source review.
\documentclass[11pt]{article}
\usepackage[margin=25mm]{geometry}
\usepackage{fontspec}
\setmainfont{Latin Modern Roman}
\usepackage{amsmath,amssymb,mathtools,unicode-math}
\setmathfont{Latin Modern Math}
\usepackage{enumitem,needspace,xurl,hyperref}
\setlength{\parindent}{0pt}
\setlength{\parskip}{5pt}
\setlength{\emergencystretch}{3em}
\newcommand{\N}{\mathbb{N}}
\newcommand{\Z}{\mathbb{Z}}
\newcommand{\Q}{\mathbb{Q}}
\newcommand{\R}{\mathbb{R}}
\newcommand{\C}{\mathbb{C}}
\newcommand{\F}{\mathbb{F}}
\newcommand{\A}{\mathbb{A}}
\newcommand{\PP}{\mathbb{P}}
\newcommand{\E}{\mathbb{E}}
\newcommand{\eps}{\varepsilon}
\newcommand{\dd}{\,\mathrm{d}}
\begin{document}
\section*{UnsolvedMath ID: 801072}
\textbf{Source problem code:} catalogue entry
\subsection{Definitions and mathematical statement}
The supplied excerpt omits the defining formula for the family \(f_{c,b}\), beginning instead with a reference to parameters \((c,b)\in\C^2\). Consequently, the following formulation is conditional on recovering the intended family; it does not identify a specific family from the excerpt alone.

For the intended complex H\'enon interpretation, suppose that, for \(b\ne0\), the missing specification gives an invertible map \(f_{c,b}:\C^2\to\C^2\). Define
\[
K^+_{c,b}=\{z\in\C^2:(f_{c,b}^{\,n}(z))_{n\ge0}\text{ is bounded}\},
\qquad
K^-_{c,b}=\{z\in\C^2:(f_{c,b}^{-n}(z))_{n\ge0}\text{ is bounded}\},
\]
and
\[
J_{c,b}=\partial K^+_{c,b}\cap\partial K^-_{c,b},
\]
where boundedness and boundaries are taken in the usual Euclidean structure on \(\C^2\). The phase space and invertibility required here must be confirmed from the omitted family definition.

The source explicitly includes \(b=0\) to simplify the presentation and defines that case separately: for \(p_c:\C\to\C\), \(p_c(z)=z^2+c\), let
\[
K(p_c)=\{z\in\C:(p_c^{\,n}(z))_{n\ge0}\text{ is bounded}\},
\qquad J_{c,0}=\partial K(p_c).
\]
Thus the definition at \(b=0\) uses the one-variable Julia set and does not require backward iteration of \(f_{c,0}\).

Define the connectedness locus
\[
\mathcal M=\{(c,b)\in\C^2:J_{c,b}\text{ is connected}\}.
\]
The target of Conjecture 24, attributed in the supplied record to Berger and Bedford (2023), is that \(\mathcal M\) is disconnected in the usual topology of the parameter space \(\C^2\). Equivalently, there exist disjoint, nonempty subsets \(U,V\subseteq\mathcal M\), both open relative to \(\mathcal M\), such that \(\mathcal M=U\cup V\). This remains an incompletely specified conjecture until the defining formula for \(f_{c,b}\) is recovered.
\subsection{Short English statement}
Conjecture 24 asserts that the set of complex parameter pairs for which the intended H\'enon Julia set is connected is itself disconnected. The slice \(b=0\) is included using the Julia set of \(z\mapsto z^2+c\). The supplied excerpt omits the H\'enon family formula, so the statement requires that missing specification.
\subsection{Sources}
\begin{itemize}
\item Local UnsolvedMath registry, supplied complete record: permanent ID \(801072\), exact source problem code \texttt{catalogue entry}, title ``Henon Maps: Open Problems --- Conjecture 24.'' The registry attributes the problem to Berger and Bedford, proposed in 2023, and labels it open; that status is supplied metadata, not an independently verified mathematical or literature assessment.
\item Berger, Bedford, \emph{Henon Maps: A List of Open Problems} (2023), Conjecture 24. Supplied source link: \url{https://arxiv.org/abs/2312.03907}. The record reports access on 2026-10-16 and extraction by \texttt{arxiv\_latex}; these are supplied provenance statements, not a claim of independent source access or verification.
\end{itemize}
\end{document}
